\documentclass{article}
\usepackage{amsmath}
\begin{document}
\section*{When retiring a reader can preserve an old transfer}
This study follows one old, keyed warehouse operation across an early
receiver retirement, delayed radio dispatch and an authority change.
The receiver being retired is a diagnostic-only participant. Its
controlled withdrawal is authorized in the fixture and does not
cancel the mandatory receipt. A distinct listener can remain bound
to the same endpoint; the policy is not authorized to withdraw it.
Only the studied simulator declarations carry these permissions.

\paragraph{Native effect and boundary.}
The broker waits for all registered participants to have a pose before
popping its queued datagrams. The test introduces one participant
without its first pose; the diagnostic and unrelated traffic are both
observably pending. After the pose arrives, it walks the endpoint's
remaining bindings. In the selected radio implementation, each
attempt that fits is debited at the sender before checking propagation.
A diagnostic directed to a distant receiver can therefore fail without
a callback and nevertheless spend bytes. Native receiver retirement
before dispatch removes that receiver's binding. Retirement afterward
does not reimburse already recorded bytes. If one other listener
remains, early retirement need not prevent an attempt.

Let $C$ be the sender's one-second allowance. The old digest delivery
has size $M=173$. At the single modeled prior-debit time $t=0.25$,
there are $d$ bytes already charged, followed by a diagnostic of $z$
bytes offered to each of $m$ remaining listeners. Write
\[
d_0=d,\qquad
d_{i+1}=\begin{cases}d_i+z&d_i+z\le C,\\d_i&\text{otherwise.}\end{cases}
\]
Under the declared good old link, fresh pose and empty old receiver
ledger, an old send through the primary at deadline $T$ fits exactly
when $M+d_m\le C$ for $T\le1.25$ (and $M\le C$ after expiry).
Equivalently $d_m=d+z\min(m,\lfloor (C-d)/z\rfloor)$ for $z>0$
and $0\le d\le C$. Retirement changes the effective $m$, not the
debit already recorded.
At $C=320,d=140,z=8,T=1.25$, an early sole binding retirement leaves
the old send possible. One remaining or previously attempted binding
leaves at least 148 live bytes, one too many. At $T=1.2501$ the budget
no longer distinguishes those runs. Negative cases keep the old
meaningful duty; an ordinary silence is never accepted as settlement.

\paragraph{Recovery scheduling.}
In the two-copy arm both pre-authorized senders already hold the same
digest. The backup also has a separate 160-byte optional job to serve.
Use the primary when its live debit permits the old send and use the
backup otherwise. The decision is local to the qualified broker and
uses the actual post-dispatch primary ledger and its known expiration,
not hidden cause or future observation. The selected envelope has no
other intervening debit; in a changed envelope both future send
windows must be rechecked. First preserve the old operation, then
maximize optional service. If the primary can serve, using it leaves
the backup able to send the optional report; if it cannot, serving
both is impossible within the two-send alternatives and fixed
capacity. The native result matches enumeration with the same ledger.
The static backup arm preserves the old duty but misses the optional
one; the static primary arm can miss the old duty. A stronger control tries the primary immediately and, if no callback
arrives, tries the backup within the same modeled tick. It also matches
old and optional service, and its charged bytes tie: the rejected
capacity attempt costs no debit in the source. It issues one more
attempt exactly when the primary cannot fit. The test does not assign
that extra invocation latency or a resource cost. There is no
demonstrated service advantage over ordinary this-tick fallback
where such feedback is implementable.

The native ingress merely queues an offered datagram. The surrounding
broker plugin normally calls dispatch at the end of the pose callback.
The test separately models a deadline with one such dispatch
opportunity and no second pose callback at the cut. That is an
assumption on the schedule, not a physical frequency bound. For this
cut an ordinary fixed hedge queues both old sends before dispatch.
At 320 bytes the backup is still charged even in runs where
the first copy has filled the receive window enough to reject a
duplicate. It delivers the old result but loses the optional report.
Sending only the backup also loses that report. A fixed schedule
that sometimes omits the backup cannot guarantee old delivery in
the indistinguishable over-budget branch. The locally informed
choice preserves the optional report whenever the primary can carry
the old data. An equally informed direct selector still matches it;
the advantage identifies a specific feedback or information
dependency, not a universal architectural win.

\paragraph{Evidence scope.}
The selected broker executes native pose-cache ingestion, queued submission, dispatch, debit,
retirement and the receiving callback. The actual old worker is
crashed after committing the old result and then fenced. The
destination has a backed checkpoint before the radio timeline. It
installs only after receiving the matching fenced digest; it returns
the existing keyed result rather than redoing the operation. Both
the unrelated radio callback and a native independent warehouse
move are checked. The radio fixture uses an authored range predicate
and controlled civil timing, not a full RF or hardware experiment.
The digest payload is fixed-length test framing; the authorization
of backups, provisioning cost, risk law, and future arrival support
are not inferred. This is a held-in recovery comparison and a
certificate distinction, not I1/I2 acceptance or an advantage over
an equally informed direct policy.
\end{document}
